\documentclass{article}

\usepackage{booktabs}
\usepackage{tabularx}
\usepackage{hyperref}

\hypersetup{
    colorlinks=true,       % false: boxed links; true: colored links
    linkcolor=red,          % color of internal links (change box color with linkbordercolor)
    citecolor=green,        % color of links to bibliography
    filecolor=magenta,      % color of file links
    urlcolor=cyan           % color of external links
}

\title{Hazard Analysis\\\progname}

\author{\authname}

\date{}

\input{../Comments.text}
\input{../Common.text}

\begin{document}

\maketitle
\thispagestyle{empty}

~\newpage

\pagenumbering{roman}

\begin{table}[hp]
\caption{Revision History} \label{TblRevisionHistory}
\begin{tabularx}{\textwidth}{llX}
\toprule
\textbf{Date} & \textbf{Developer(s)} & \textbf{Change}\\
\midrule
Date1 & Name(s) & Description of changes\\
Date2 & Name(s) & Description of changes\\
... & ... & ...\\
\bottomrule
\end{tabularx}
\end{table}

~\newpage

\tableofcontents

~\newpage

\pagenumbering{arabic}

\wss{You are free to modify this template.}

\section{Introduction}

\wss{You can include your definition of what a hazard is here.}

\section{Scope and Purpose of Hazard Analysis}

\wss{You should say what \textbf{loss} could be incurred because of the
hazards.}

\section{System Boundaries and Components}

This section divides DAIIM into its major components and identifies where the
boundary of the system lies. Components inside the boundary are built,
deployed, and maintained by the team, and their behaviour can be changed to
mitigate hazards. Components outside the boundary are part of DAIIM's
environment: DAIIM depends on them, but cannot control them, so hazards
originating from them can only be detected, tolerated, or contained. The
components below are based on the requirements in the SRS and the technologies
listed in the Development Plan. They are preliminary and will be refined in the
design documents.

\subsection{Components Within the System Boundary}

\begin{enumerate}

\item \textbf{Web Client.} The browser-based user interface, built with React,
through which task requesters and inference providers use DAIIM. It connects to
the user's wallet for authentication, lets users browse and create tasks,
submit inference requests, register as providers, and view results,
reputations, stakes, and rewards. Before each transaction, it shows the user
the action being performed and the amount of tokens being transferred, and
afterwards it shows the transaction's status and identifier. It also performs
an initial check that request inputs conform to the task's input schema.
\textit{Related SRS requirements:} Accounts and Authorization, Inference Tasks,
Inference Requests, Aggregation and Results, Transparency.

\item \textbf{On-Chain Programs.} The Solana programs, written in Rust, that
hold DAIIM's authoritative state and all of its tokens. They are responsible
for:
\begin{itemize}
  \item the registry of inference tasks and their parameters (minimum
  providers, reputation tolerance, acceptance period, maximum prediction time);
  \item provider registration and the escrow of provider stakes;
  \item escrow of requester payments and the lifecycle of each inference
  request (acceptance period, prediction phase, cancellation and refunds);
  \item recording predictions, aggregating them into the final result, and
  updating provider reputations;
  \item distributing rewards, and slashing or burning provider stakes.
\end{itemize}
Because this component holds funds and blockchain transactions cannot be
reversed, a defect here can cause direct and permanent financial loss.
\textit{Related SRS requirements:} Inference Tasks, Provider Registration and
Staking, Provider Acceptance and Predictions, Aggregation and Results,
Reputation and Rewards, Penalties, Transparency.

\item \textbf{Off-Chain Backend Service.} A Node.js and TypeScript service that
performs the work that is impractical or unsafe to do on-chain. It monitors the
on-chain programs for new inference requests and broadcasts them to the
providers registered for the task, stores request input data and distributes it
only to the requester and the providers that accepted the request, and deletes
the input data once the request is completed or cancelled. It may also submit
the transactions that move a request from one phase to the next when a deadline
passes, since on-chain programs cannot act on their own, and it serves
marketplace information to the Web Client. Since this service is a single
point that the team operates, it must not be able to change results, payments,
or reputations, which are controlled only by the on-chain programs.
\textit{Related SRS requirements:} Inference Requests, Provider Acceptance and
Predictions, Data Handling.

\item \textbf{Off-Chain Database.} A PostgreSQL database used by the backend
service. It temporarily holds encrypted request input data, which must not be
stored on the blockchain, as well as task information that is too large to
store on-chain economically (such as problem descriptions, schemas, and
references to training data) and indexed copies of on-chain records for fast
browsing. The on-chain records remain the source of truth for any data that
appears in both places. \textit{Related SRS requirements:} Inference Tasks,
Data Handling.

\item \textbf{Provider Interface.} The interface, such as an API or client
library, through which an inference provider's model receives broadcast
requests, accepts requests, retrieves input data, and submits predictions. It
checks that predictions conform to the task's output schema before they are
submitted. This component ends where the provider's own model begins.
\textit{Related SRS requirements:} Provider Acceptance and Predictions.

\end{enumerate}

\subsection{Components Outside the System Boundary}

\begin{enumerate}

\item \textbf{Solana Blockchain Network.} The public Solana network, including
its validators and the RPC nodes used to read state and submit transactions.
DAIIM relies on it to execute the on-chain programs, store records, and provide
the time used for request deadlines. It may become congested, raise fees, delay
or drop transactions, or suffer outages.

\item \textbf{User Wallets.} Third-party Web3 wallets, such as browser wallet
extensions, that hold users' private keys and sign transactions. DAIIM cannot
recover a lost key, and cannot distinguish a compromised wallet from its
legitimate owner.

\item \textbf{Inference Providers' Models and Infrastructure.} The machine
learning models, computing resources, and network connections of each
provider. DAIIM receives only the predictions they produce and cannot inspect
the models. Providers may be unavailable, inaccurate, careless, or dishonest,
and several providers may collude.

\item \textbf{Task Requesters and Their Input Data.} The people who define tasks
and submit requests, and the sources of the input data they provide. DAIIM can
check that inputs conform to a task's schema, but cannot verify that the data
is correct, that task parameters are sensible, or that the data is not
sensitive.

\item \textbf{Downstream Decisions.} The decisions that task requesters and
beneficiaries, such as emergency management agencies, make using DAIIM's
results. DAIIM is responsible for providing accurate results and information
about their quality, but not for how the results are used.

\item \textbf{Hosting Infrastructure and Browsers.} The cloud or server
infrastructure that hosts the backend service and database, and the web
browsers on which the Web Client runs.

\end{enumerate}

\section{Critical Assumptions}

\wss{These assumptions that are made about the software or system.  You should
minimize the number of assumptions that remove potential hazards.  For instance,
you could assume a part will never fail, but it is generally better to include
this potential failure mode.}

\section{Failure Mode and Effect Analysis}

\wss{Include your FMEA table here. This is the most important part of this document.}
\wss{The safety requirements in the table do not have to have the prefix SR.
The most important thing is to show traceability to your SRS. You might trace to
requirements you have already written, or you might need to add new
requirements.}
\wss{If no safety requirement can be devised, other mitigation strategies can be
entered in the table, including strategies involving providing additional
documentation, and/or test cases.}

\section{Safety and Security Requirements}

\wss{Newly discovered requirements.  These should also be added to the SRS.  (A
rationale design process how and why to fake it.)}

\section{Roadmap}

\wss{Which safety requirements will be implemented as part of the capstone timeline?
Which requirements will be implemented in the future?}

\newpage{}

\section{Appendix --- Generative AI Use Disclosure}

This section provides an overview of how generative AI tools were used in the
making of this document.

\subsection{AI Tools Used}

\begin{itemize}
    \item Claude Opus 5.5 (Anthropic), used through Claude Code.
\end{itemize}

\subsection{How and Where Used}

Claude Opus 5.5 was used throughout this document to help ideate and draft its
content. The team provided the tool with the SRS and the Development Plan, and
used it to brainstorm the hazards, components, assumptions, and requirements
that the document should address, and to produce first drafts of the text.
In Section~3 (System Boundaries and Components), for example, it proposed a
division of DAIIM into components inside and outside the system boundary,
drafted the description of each component, and linked each component to the
groups of functional requirements in the SRS that it is responsible for. It was
also used to draft the text of this appendix.

The tool was used as an aid, not as a replacement for the team's own work.
Team members were involved in every sentence of the document: each AI-assisted
draft was read, discussed, and edited by the team, and ideas that did not fit
the project were discarded or rewritten.

\subsection{Verification of AI-Generated Content}

Every sentence in this document was reviewed by at least one team member
before it was included. AI-assisted content was checked against the SRS and
the Development Plan to confirm that it describes responsibilities and
behaviour that the requirements actually assign to DAIIM, that it does not
introduce requirements or design decisions the team has not agreed to, and
that it is consistent with the technologies the Development Plan lists.
References to SRS requirements were checked by hand. Content that is still
undecided, such as how request phases are advanced when deadlines pass, is
marked as preliminary and will be confirmed in the design documents. The team
remains responsible for the accuracy, completeness, and consistency of this
document, regardless of whether material was first drafted by a team member or
with the help of an AI tool.

\newpage{}

\section{Appendix --- Reflection}

\wss{Not required for CAS 741}

\input{../Reflection.text}

\begin{enumerate}
    \item What went well while writing this deliverable? 
    \item What pain points did you experience during this deliverable, and how
    did you resolve them?
    \item Which of your listed risks had your team thought of before this
    deliverable, and which did you think of while doing this deliverable? For
    the latter ones (ones you thought of while doing the Hazard Analysis), how
    did they come about?
    \item Other than the risk of physical harm (some projects may not have any
    appreciable risks of this form), list at least 2 other types of risk in
    software products. Why are they important to consider?
\end{enumerate}

\end{document}